\section{Day 3: (Sep.\ 15, 2026)}
It turns out that Hunter Spink types out his lecture notes! Regardless, I'll try to keep up here (until I burn out).

Recall that we defined $\abs{\Spec R}$ to be the collection of prime ideals in $R$. For example, $\abs{\Spec \RR[x]}$ is given by $\CC$ modded out by complex conjugation, along with the $\{0\}$ ideal itself (we call this the transcendental solution to $Z(\emptyset)$).

Recall that $\Hom_{\RR-\opname{alg}}(C^0(M, \RR), \RR) = \{\ev_p \colon p \in M\}$. The idea being, if you want to identify $Z(f)$, we can simply look at all the points $p$ such that $\ev_p f = 0$ instead. Specifically, we want to look at the bijection between $\{\ev_p \colon p \in M\}$ and $\{p \in M\}$.

Consider $\abs{\Spec K}$, where $K$ is a field. Clearly, $\abs{\Spec K} = \{(0)\}$. Using the fact that $K(p) = \opname{Frac}(R/p)$ and $f \in R$, we have that $f(p)$ is the inverse of $f$ in $K(p)$, i.e., $\ev_p \colon R \to K(p)$.

\begin{definition}
    For $\SA \subset R$, we define $Z(\SA)$ to be
    \[ Z(\SA) = \{p \in \Spec R \mid f(p) = 0, \forall f \in \SA\} \subset \Spec R. \]
    Specifically, the above definition is equivalent to $p \in \Spec R$ such that $\SA \subset p$. Think of the above as ``all functions that vanish at $p$''.
\end{definition}
We define the Zariski topology on $\abs{\Spec R}$ to have $Z(\SA)$ as the closed sets. This is a \textit{tentative definition}, and will be expanded upon later. We want to show that this is actually a topology!

\begin{theorem} \label{thm:Zariski_topology}
    The Zariski topology (described above) is a topology.
\end{theorem}
\begin{proof}
    First, we want to show that $\emptyset$ and $\abs{\Spec R}$ lie in the Zariski topology. For $\emptyset$, we can take $Z(R)$, or really just $Z(1)$ ($1$ being picked because it is a unit). $1$ has the great feature that it never vanishes in any field. For $\abs{\Spec R}$, we can take $Z(0) = Z(\emptyset)$.

    For finite unions (since we're prescribing closed, not open sets), we have that $Z(\SA) \cup Z(\SB) = Z(\SA \SB)$. For arbitrary interscetions, we have
    \[ \bigcap Z(A_\beta) = z\left(\bigcup A_\beta\right). \]
    Now, we want to check that
    \[ Z(\SA \SB) = \{p \mid fg \in p, \forall f \in \SA, g \in \SB\}, \]
    and
    \[ Z(\SA) \cup Z(\SB) = \{p \mid f \in p, \forall f \in \SA\} \cup \{p \mid g \in p, \forall g \in \SB\}, \]
    in order to check that our statement for finite unions holds; indeed, this follows from the fact that, if there exists $f \in \SA$ with $f \notin p$ and $g \in B$ with $g \notin p$, then $fg \notin p$.
\end{proof}

We want to look at some examples of topological spaces now. As an example, consider $\Spec \ZZ$. This space has one point for every prime ideal in $\ZZ$, of which are the zero ideal and ideals generated by primes. What does $Z(4)$ look like here? $\ol 4 = 0$ in $(2) = \FF_2$, $\ol 4 = 1$ in $(3) = \FF_3$, and so on, with it simply being $4$ in $(0) = \ZZ$. Thus, $Z(4) = \{(2)\}$ (i.e., check the prime factorization).\footnote{four, four, four, four... (draws a lot of $4$s on the blackboard). JHIN!!!!!} We can try this with other numbers too, such as $Z(12) = \{(2), (3)\}$, $Z(1) = \emptyset$, and $Z(0) = \Spec \ZZ$.

The closed sets here are $\emptyset$, $\Spec \ZZ$, and finite collections of primes. In particular, $\ol{\{(0)\}} = \Spec \ZZ$. At this point, Hunter made a remark that the facts we know about manifold topology don't apply here at all, which makes sense, since this space is... really weird.

Recall that
\[ \abs{\Spec \RR[x]} = \left(\frac{\CC}{z \leftrightarrow \ol z}\right) \cup \{(0)\}. \]
Here, $Z(f) = Z(f_1) \cup \dots \cup Z(f_k)$, given that we can factor $f$ as $f_1 f_2 \dots f_k$ with each $f_i$ irreducible. For irreducible $f$, if $f = x - a$, then $Z(f) = \{(x - a)\}$; if $f = (x - a)^2 + b^2$, then $p = (f)$. Both $\ZZ$ and $\RR[x]$ are UFDs; in particular, they are Dedekind domains, i.e., all nonzero proper ideals can be factored into a product of prime ideals.

We're now going to do an example where the ring is not a UFD. Consider $\abs{\Spec \CC[x, y]}$ (where $\CC[x, y]$ is a UFD); per the Nullstellensatz, there are three types of primes. We can think of the dimension $0$, $1$, and $2$ varieties, which will correspond to these prime ideals. We classify by height of primes, i.e., the transcendence degree of $\kappa(p)/\CC$. For dimension $0$ varieties, we can pick out individual points $(a, b)$ of the affine space $\mathbb A^2$. For $f$ irreducible, $\ol{\{(f)\}} = (f) \cup \{(x - a, y - b) \mid f(a, b) = 0\}$. ``As a topological space, this thing is a disaster.''

When does $Z(g)$ contain $(f)$? This occurs when $g = 0$ in $\opname{Frac}(\CC[x, y] / (f))$, which is equivalent to $g \in (f)$. Then $g$ is a multiple of $f$. Since $\CC[x, y]$ is a UFD, the smallest $Z(s)$ contains $f$ if $Z(f)$ is given by $p$ such that $f(p) = 0$, i.e., $\{(x - a, y - b), f(x, y) = 0, (f)\}$. You can always find closures from first principles, but to be more efficient you need ``non-tautologous math''.